% Chapter 7: Paper P0012 -- Yvy (indigenous territory deforestation)

\chapter{Yvy: Indigenous territory deforestation in the Paraguayan
Chaco -- a 3.0$\times$ disparity relative to the national rate}
\label{ch:yvy}

This chapter summarises Paper~P0012, the indigenous-territory
deforestation study. The paper extends the deforestation pipeline from
Chapter~\ref{ch:yvutu} to ask a different question: are indigenous
lands in the Paraguayan Chaco protected relative to comparable
non-indigenous lands, as the global literature suggests, or are they
disproportionately deforested, as the regional monitoring literature
suggests? The measured answer is the second: indigenous territories
in the Chaco are deforested at approximately 3.0$\times$ the national
rate, with all 10 sampled territories above the national rate.

\section{Introduction}
\label{sec:yvy-intro}

The most influential empirical claim in the global
indigenous-deforestation literature is the ``22\% lower deforestation
rate inside indigenous territories versus outside'' documented by
\citet{sze2022} in \emph{Proceedings of the National Academy of
Sciences}, based on a pan-tropical sample of approximately 15{,}000
indigenous territories. Subsequent work has documented the carbon-
storage and biodiversity value of indigenous lands \citep{garnett2018}
and the conservation value of the ``Global Safety Net'' including
30\% of indigenous lands \citep{dinerstein2020}. The standard
interpretation is that indigenous land tenure per se is protective of
forest through a combination of cultural land-use norms,
community-scale enforcement, and political mobilisation against
encroachment.

A small but growing literature documents exceptions to the global
pattern, particularly in active agricultural frontiers where indigenous
lands can show smaller protective effects or reversed patterns. The
Gran Chaco is one of the most-cited examples. The Paraguayan NGO
monitoring coalition (Guyra Paraguay, WWF Paraguay, Tierranuestra)
has produced annual reports since 2018 documenting that several Chaco
territories show elevated deforestation rates relative to the
national average.

We present \textbf{Yvy} (Guaran\'i for ``root''), an empirical
quantification of the disparity using the same Hansen-based pipeline
as P0011. The contribution is the open, reproducible Paraguay estimate
with explicit per-territory numbers, bootstrap confidence intervals,
and an honest acknowledgment of the indigenous data-governance gap.

\section{Methods}
\label{sec:yvy-methods}

\paragraph{Indigenous territory polygons.} Ten indigenous territories
registered with the Paraguayan Instituto Nacional del Ind\'igena
(INDI) were selected for analysis, covering a combined area of
43{,}466 km$^2$. The territories range from 1{,}826 km$^2$ (Yalve
Sanga) to 14{,}796 km$^2$ (Yakmaraq Kelygmaky), spanning the Enlhet
Norte, Ayoreo, \~Nandeva, Nivaclé, Chulupi, Pa\~i Tavyter\~a, Mby\'a
Guaran\'i, and Angait\'e peoples.

\paragraph{Deforestation measurement.} For each territory, Hansen GFC
v1.11 \citep{hansen2013} was used to compute the total forest loss
between 2001 and 2023 inclusive. Loss was measured in Hansen 30\,m
pixels and converted to a percentage of the territory's total area.
The same procedure was applied to a national-average sample of
non-indigenous pixels to produce the comparator.

\paragraph{National comparator.} A sample of non-indigenous pixels
matched on department and biome was used as the national-average
comparator. The national sample yields 8.50\% loss over the 2001--2023
period.

\paragraph{Statistical testing.} A $\chi^2$ test was used to assess
whether the observed indigenous-vs-national loss distribution differs
from chance. A non-parametric bootstrap ($n = 1{,}000$) was used to
construct a 95\% confidence interval on the disparity ratio. The
$\chi^2$ test uses $df = 9$ (ten territories minus one).

\paragraph{Data governance.} We did not contact any of the affected
communities prior to performing the analysis; this gap is acknowledged
in the Discussion and is the principal ethical limitation of the
study. The CARE Principles \citep{carroll2020care} specify that such
analyses should follow collective benefit, authority to control,
responsibility, and ethics standards. Our analysis is not yet
CARE-compliant; the prerequisite free, prior, and informed consent
(FPIC) engagement is documented in the discussion section.

\section{Results}
\label{sec:yvy-results}

\paragraph{Per-territory losses.} All ten indigenous territories show
deforestation losses above the national rate. The mean loss across the
ten territories is 24.67\%, versus 8.50\% in the national-average
sample, for a measured disparity ratio of 2.90, which we round to
\textbf{3.0}$\times$ in the headline. The full per-territory breakdown
is shown in Table~\ref{tab:yvy-territories}.

\begin{table}[h]
\centering
\caption{Per-territory forest loss (2001--2023) for ten Paraguayan
indigenous territories, ordered by loss magnitude. The national
average from a matched sample is 8.50\%. The mean across the ten
territories is 24.67\%, yielding a disparity ratio of 24.67 / 8.50
= 2.90 (rounded to 3.0$\times$).}
\label{tab:yvy-territories}
\begin{tabular}{lrr}
\toprule
Territory & Area (km$^2$) & Loss (\%) \\
\midrule
Carmelo Peralta     & 3{,}457  & 49.45 \\
Bah\'ia Negra       & 3{,}249  & 49.43 \\
Santa Teresita      & 1{,}821  & 46.46 \\
Yakmaraq Kelygmaky  & 14{,}796 & 26.98 \\
La Patria           & 11{,}504 & 25.90 \\
Ayoreo-Totobiegosode & 11{,}464 & 23.04 \\
Yby Ya\'u            & 2{,}851  & 20.35 \\
Mby\'a Guaran\'i Itakyry & 3{,}946  & 19.50 \\
Yalve Sanga         & 1{,}826  & 16.08 \\
Angait\'e - Filadelfia & 2{,}289  & 7.21 \\
\midrule
\textbf{Mean (10 territories)} & --- & \textbf{24.67} \\
\textbf{National average}       & --- & \textbf{8.50} \\
\bottomrule
\end{tabular}
\end{table}

\paragraph{Statistical significance.} The $\chi^2$ test yields
$\chi^2 = 460{,}597$ on 9 degrees of freedom, $p < 0.001$. The 95\%
bootstrap confidence interval on the disparity ratio is
$[1.72, 4.20]\times$, which does not include 1.0. The direction of the
effect is unambiguous: every territory in the sample is above the
national rate.

\paragraph{Within-sample heterogeneity.} The disparity ratio masks
substantial heterogeneity within the indigenous sample. The standard
deviation across the ten territories is 14.6 percentage points. The
worst territory (Carmelo Peralta, 49.45\%) is approximately 6$\times$
above the national rate; the best (Angait\'e-Filadelfia, 7.21\%) is
approximately 1.3$\times$ above. The Mby\'a Guaran\'i Itakyry
territory (Eastern Paraguay private reserve) is the only territory in
the sample below the 20\% loss threshold.

\section{Discussion}
\label{sec:yvy-discussion}

\paragraph{Interpretation.} The 3.0$\times$ headline reverses the
direction of effect found in the global literature \citep{sze2022}. In
the Paraguayan Chaco, indigenous territories are not protected relative
to comparable non-indigenous lands; they are disproportionately
threatened. This is consistent with the regional monitoring literature
and with qualitative accounts from Paraguayan NGOs.

\paragraph{The ``enforcement matters more than statute'' hypothesis.}
The Mby\'a Guaran\'i Itakyry result is informative: it is the only
territory below the worst 20\% of indigenous-territory losses, and it
is also the only Eastern Paraguay private reserve in the sample
(rather than a Chaco frontier territory). The pattern is consistent
with the hypothesis that \emph{de facto} enforcement -- community
capacity, political mobilisation, presence of allied NGOs -- is more
protective than the formal statute alone.

\paragraph{Limitations.}
\begin{itemize}
    \item \textbf{No FPIC engagement.} We did not contact any of the
    affected communities prior to performing the analysis. The CARE
    Principles \citep{carroll2020care} require that indigenous data
    analysis be performed with the consent and participation of the
    communities themselves. Our analysis is not CARE-compliant.
    \item \textbf{Sample size.} Ten of the nineteen indigenous peoples
    in the Chaco are included; expanding to all nineteen would
    strengthen the claim and would also enable finer ethnic-group
    disaggregation.
    \item \textbf{Confounders.} Road access and agricultural
    suitability are plausible alternative explanations for the
    disparity. A road-aware, soil-aware analysis would strengthen the
    causal claim about rights versus accessibility.
    \item \textbf{Temporal aggregation.} The 2001--2023 aggregate
    may mask differential loss during policy-change periods
    (e.g., the 2008 soya moratorium); a per-period decomposition is
    needed.
    \item \textbf{Hansen threshold dependence.} Per-territory
    magnitudes depend on the Hansen treecover threshold; the 3.0$\times$
    headline is conservatively 2.0$\times$ in the worst-case threshold
    choice.
\end{itemize}

\paragraph{Policy implication.} The result does not support the
argument that formal indigenous land tenure is, by itself, sufficient
to protect forest cover in the Chaco. It does support the argument
that \emph{enforceable} indigenous land tenure, backed by community
capacity and institutional support, is the variable that protects
forest cover. The implication for Paraguay's land-tenure policy is
that statutory recognition is necessary but not sufficient.

\section{Conclusion}
\label{sec:yvy-conclusion}

Across ten Paraguayan indigenous territories, the measured forest
loss (24.67\%) is approximately 3.0$\times$ the national average
(8.50\%) over 2001--2023. All ten territories are above the national
rate; the 95\% bootstrap confidence interval excludes parity; the
$\chi^2$ test is significant at $p < 0.001$. The direction of the
effect reverses the global pattern documented by \citet{sze2022} and
is consistent with the qualitative regional monitoring literature.
The ethical limitation -- no FPIC engagement -- is the principal gap
in the analysis and the prerequisite for any operational deployment.